\chapter{MIPS/TS: a terminal session}\label{mips-ts}
\index{MIPS/TS}\index{SOS}\index{TEM}\index{PRIMO}

\section{What this environment provides}\label{mips-ts-overview}

MIPS/TS adds terminal and job services to the monitor. The supplied profile
starts three resident subsystems: TEM manages terminal sessions, SOS handles
account login and the user's process, and PRIMO provides peripheral transport
services. SOS is the command interface used to log in; TEM is entered by the
\texttt{run} command. Once enrolled, the user works in FP, using the guest
commands described elsewhere in this manual. SOS is not a replacement for
either FP or the emulator's \texttt{sim4000>} operator interface.

Historically, SOS means \emph{Swopping Online System}. A SOS job is a
process created from a named description in its user catalog. In this
demonstration, \texttt{test} names that description: it supplies the initial
program, memory limits, password and disk claims. It is not a host account,
nor does logging in create a Unix-like home directory. The FP command
\texttt{scope user} describes guest catalog retention, not a host user's
ownership or a security boundary.

This is a ready-to-try demonstration with one configured account. Automatic
startup, TEM-mediated login, incorrect-password rejection, FP commands and
normal logout have been checked. It uses the compatibility monitor with IX
enabled and storage protection off. Keep this unprotected profile for MIPS/TS
operation. The standalone protected-monitor profile is separate; do not enable
protection in this configuration or substitute its disk images.

\begin{operationnote}{Local demonstration credentials only}
The account and password are both \texttt{test}. Telnet has no encryption;
the supplied listener binds only to \texttt{127.0.0.1}. Do not expose this
profile to a network or treat it as a secure multiuser service.
\end{operationnote}

\subsection{Choose a checked workflow}\label{mips-ts-workflow-guide}

For a first visit, start the supplied profile, log in, try
\texttt{message}/\texttt{lookup} and finish normally. The later exercises are
independent: printing does not require the retained-file or batch exercise.
Use one \texttt{test} job at a time, and finish it before submitting batch work.

\begingroup
\footnotesize
\begin{longtable}{@{}>{\raggedright\arraybackslash}p{0.23\linewidth}>{\raggedright\arraybackslash}p{0.28\linewidth}>{\raggedright\arraybackslash}p{0.41\linewidth}@{}}
\toprule
Task & Main command context & Procedure \\
\midrule
\endhead
\bottomrule
\endfoot
First terminal session & Host shell, then remote SOS/FP &
\manualref{mips-ts-start}{Start}, then \manualref{mips-ts-login}{log in} \\
Keep a small guest file & Preparation FP, then remote FP &
\manualref{mips-ts-retained-file}{Retained-file exercise} \\
Run a small jobfile & Preparation FP, then local SOS &
\manualref{mips-ts-batch-jobs}{Batch exercise} \\
Print guest text & Local PRIMO and administrative FP &
\manualref{mips-ts-primo-printing}{Sequential printing} \\
Resume a lost connection & Remote FP; no new login &
\manualref{mips-ts-idle-reconnect}{Reconnect}, including \manualref{mips-ts-partial-input}{known partial input} \\
Check a job or service & Local guest system \texttt{s} &
\manualref{mips-ts-status}{Process status} \\
Unexpected failure & Local operator; preserve the job &
\manualref{mips-ts-unexpected-failure}{Preserve evidence and work} \\
Stop or back up & Remote FP, then local operator &
\manualref{mips-ts-normal-stop}{Normal stop and backup} \\
\end{longtable}
\endgroup

Expected diagnostics are stated beside the relevant commands. In particular,
an intentional \texttt{lookup ... unknown} is different from a failed import,
an abnormal SOS removal or an unfinished PRIMO transport. Do not continue past
an unexpected error just because another prompt or bell appears.

\section{Select the files and start}\label{mips-ts-start}

The ordinary build directory and binary archive contain the prepared files:

\begin{center}
\begin{tabular}{@{}ll@{}}
\toprule
File & Purpose \\
\midrule
\texttt{sim4000} & Emulator executable \\
\texttt{MONITOR} & Compatibility monitor, shared with the default profile \\
\texttt{UNIT0.mipsts} & System disk, account catalog and automatic startup \\
\texttt{UNIT1.mipsts} & Prepared second disk, mounted as \texttt{disc1} \\
\texttt{RC4000.mipsts.cfg} & Matching devices, Telnet and persistence settings \\
\bottomrule
\end{tabular}
\end{center}

Keep these files together and create \texttt{exchange/} for host transfers.
No Python, preparation utility or guest rebuilding is needed to run them.
Do not substitute the default or protected disk for either MIPS/TS disk.

In a host shell, change to the build or unpacked distribution directory:

\begin{commandexample}[language=sh]{Host shell: start the MIPS/TS profile}
./sim4000 -a -c RC4000.mipsts.cfg
\end{commandexample}

The disk's startup file mounts the second disk, sets a fixed test date of
30 September 1986, and starts TEM, PRIMO and SOS without a date prompt.
Wait for \texttt{sos started}; the release messages can arrive in a different
order. A banner alone is not a login prompt. Leave this terminal running;
it remains the local monitor/operator console.

\section{Connect and log in}\label{mips-ts-login}

Open a second host terminal and connect to the configured port:

\begin{commandexample}[language=sh]{Second host terminal: connect}
telnet 127.0.0.1 4004
\end{commandexample}

Press Ctrl-A to request guest attention. After \texttt{att}, type
\texttt{sos} and press Enter. At the SOS prompt enter the complete line:

\begin{commandexample}{SOS: start a session through TEM}
run test pass test
\end{commandexample}

A successful enrollment reports \texttt{from tem} and
\texttt{test enrolled}, followed by the FP input prompt (a bell rather than
a printed shell-style prompt). The inline password syntax above is the
checked procedure; do not rely on a separate interactive password dialogue.

\begin{commandexample}{Remote FP: a first session}
message hello
lookup fp
finis
\end{commandexample}

The first command prints a message; the second identifies the FP program
on the system disk. \texttt{finis} ends the user process, and SOS reports
\texttt{test removed after finis}. To log in again, request attention with
Ctrl-A, select \texttt{sos}, and repeat the login command. Only one terminal
account is provisioned; adding a host listener alone does not configure an
additional SOS terminal or account.

\subsection{Recognize the command context}\label{mips-ts-command-contexts}

The same terminal carries commands for different interpreters. Attention
selects a recipient; it does not log out, remove a job or suspend the whole
emulator. Enter one command at a time and wait for its complete reply.

\begin{longtable}{@{}p{0.23\linewidth}p{0.32\linewidth}p{0.37\linewidth}@{}}
\toprule
Context & How to recognize or select it & Commands in this chapter \\
\midrule
\endhead
\bottomrule
\endfoot
Host shell & Before starting the emulator, or a separate host terminal &
\texttt{./sim4000}, \texttt{telnet} \\
Emulator operator & Ctrl-E in the original terminal; \texttt{sim4000>} &
\texttt{continue}, \texttt{quit} \\
Guest system \texttt{s} & Escape in the original terminal, then \texttt{s} &
\texttt{list} for process inspection \\
SOS & Ctrl-A in Telnet, then \texttt{sos}; selection ends at \texttt{>} and a bell &
\texttt{run test pass test} or \texttt{go test pass test} \\
FP in job \texttt{test} & Successful enrollment, then the FP input bell &
\texttt{message}, \texttt{lookup}, \texttt{finis} \\
\end{longtable}

Do not type \texttt{run test pass test} at \texttt{sim4000>}: the emulator's
\texttt{run} command controls CPU execution, not a SOS login. Similarly,
\texttt{finis} belongs to the user FP, whereas \texttt{quit} exits the host
emulator. TEM's \texttt{from tem} header identifies the terminal path; it
does not mean that ordinary FP commands must be entered at a TEM setup prompt.
Ctrl-E in the original terminal suspends the entire emulated machine,
including remote jobs. Enter \texttt{continue} at \texttt{sim4000>} before
expecting more Telnet output. Guest attention alone does not suspend CPU
execution.

\subsection{Direct and TEM-mediated sessions}\label{mips-ts-go-and-run}

The recommended login uses \texttt{run}, routing terminal traffic through
SOS and TEM. TEM supplies terminal-output spooling. A direct alternative,
also checked with this demonstration account, is:

\begin{commandexample}{SOS: alternative direct session, after finishing the previous job}
go test pass test
\end{commandexample}

\texttt{go} uses SOS without TEM terminal spooling. It still creates the
same kind of FP job with the account's resource limits; it does not switch
monitor profiles or disable a device. Expect \texttt{test enrolled} but no
\texttt{from tem} header in the direct-login reply. Use one form or the
other, not both for simultaneous logins to this single-account setup.

\section{Small jobs and resource limits}\label{mips-ts-small-jobs}
\index{MIPS/TS!resource claims}

The first \texttt{message}/\texttt{lookup} exercise is intentionally small.
The login starts FP, not a complete modern shell environment; guest program
names still refer to installed RC catalog entries. See
\manualref{files-catalog-entries-and-algol-programs}{Files, catalog entries,
and Algol programs} for the general FP notation, but do not copy its process
creation or resource assumptions into a SOS session.

The demonstration description permits roughly 20,000--30,000 RC halfwords
of job memory. Small FP commands and fixed batch allocation/copy examples
have been checked; the larger Algol tutorial uses a different process with
larger claims. A resource error here is not by itself an emulator defect.

\begin{operationnote}{Small input does not imply a small allocation}
The supplied \texttt{fromhost} can allocate 36 segments even for a short
text. That exceeds this account's 30-segment temporary claim. Do not use
the general host-import tutorial unchanged in the demo job, and do not
increase free disk capacity or relabel \texttt{disc1} to address an account
claim error. The account's allowance, not merely the length of the host text,
limits allocation; device access is a separate requirement.

An enlarged account alone does not grant access to the host device.
Without explicit device inclusion, a larger-account check receives
\texttt{reserve host-file process result 2} from a SOS job. Treat that message
as failure even though it has no leading asterisks; an empty guest file may
already have been created. Use the administrative staging procedure for the
supplied demo account.

\end{operationnote}

\subsubsection{Beyond the supplied account}\label{mips-ts-provisioned-accounts}

The supplied images contain only \texttt{test}; editing a host configuration
does not create guest accounts or enlarge their claims. A separately
provisioned larger account has been checked for Algol compilation and retained
program execution across restart. It is not an additional ready-made profile.

For that larger account, direct transfers have been
checked after the operator enters \texttt{include 63} at the \emph{SOS}
command prompt, after job enrollment. Expect a timestamped
\texttt{sos: ready}, not an FP bell. This grants host-device access to
\emph{all currently enrolled jobs}; it does not enlarge disk claims, configure
future logins, or provide separate host directories per account. Only one job
can reserve the device at a time. Repeat the deliberate inclusion after
enrollment following a restart. Do not add this command to the demo's normal
startup or treat it as permission to use the host directory indiscriminately.
Two separately provisioned accounts also pass simultaneous sessions on devices
4 and 5, with separate standard/user
bases and two SOS child-process slots. An extra Telnet listener alone does
not provide this setup: equal standard bases cause \texttt{user conflict},
and insufficient child-process claims prevent the second job's creation.
Normal logout and relogin are checked for that two-account setup. The
\manualref{mips-ts-idle-reconnect}{idle reconnect procedure} below instead
covers the supplied single-account environment. Do not add an account or
replace \texttt{soscat} in a working installation merely to follow a recipe
for the supplied \texttt{test} account.

The corresponding checked batch fixtures import text in an administrative
preparation process, copy it to a compact area, then make the compact file
visible to the job. This is not an automatic feature of a remote import.
Creating more accounts or changing \texttt{soscat} requires a deliberate
guest-administration workflow; editing \texttt{RC4000.mipsts.cfg} alone
does not change SOS memory or disk claims.

\subsection{Background jobs and jobfiles}\label{mips-ts-batch-jobs}
\index{MIPS/TS!batch jobs}

SOS also offers background execution using \texttt{batch}. A jobfile is an
existing retained guest text file containing the FP commands to execute;
it is not a host shell script or a filename beneath \texttt{exchange/}.
The following small exercise has been checked on the supplied unprotected
profile, including both password forms and cleanup. Finish any interactive
\texttt{test} session first: this recipe does not demonstrate concurrent jobs.

\subsubsection{Prepare a compact jobfile}\label{mips-ts-batch-prepare}

Use private runtime copies and their overlays. The supplied image has no
\texttt{batchjob}; do not overwrite an entry left by another exercise.
Create a new host file without replacing an existing one:

\begin{commandexample}[language=sh]{Host shell: create the batch input}
(set -C; printf 'message batchprobe\nfinis\n' > exchange/batchinput)
\end{commandexample}

Its two lines are FP commands, not SOS commands:

% mips-batch: input
\begin{commandexample}{Host file contents: exchange/batchinput}
message batchprobe
finis
\end{commandexample}

In the original terminal press Escape, select \texttt{s} and allocate an
administrative preparation FP with intervals matching the demo account:

% mips-batch: administrator
\begin{commandexample}{Local Guest-s: allocate batch preparation}
all bprep base 0 0 login 0 0 user 0 0 project 0 0 run
\end{commandexample}

Wait for \texttt{to bprep} and the FP bell, then check the name:

% mips-batch: absent
\begin{commandexample}{Local FP bprep: check the target is absent}
lookup batchjob
\end{commandexample}

Expect \texttt{***lookup batchjob unknown}. If the name is found, stop and
preserve that file. Otherwise prepare the compact retained guest input:

% mips-batch: prepare
\begin{commandexample}{Local FP bprep: import and retain the jobfile}
batchtmp=fromhost batchinput
batchjob=copy batchtmp
scope user batchjob
lookup batchjob
finis
\end{commandexample}

The descriptor should say \texttt{=set 1 disc} and \texttt{; user}. Wait for
\texttt{pause bprep finis fp}. The import uses administrative resources and
\texttt{copy} compacts it; the batch job must not try to import its own input
with the demo account's smaller temporary claim. Retention keeps the jobfile
available after the preparation FP finishes. The 0/0 intervals establish
visibility; a host pathname by itself does not.

\subsubsection{Submit and recognize normal completion}\label{mips-ts-batch-submit}
\index{SOS!batch jobs}\index{SOS!password forms}

Still in the original terminal, press Escape, select \texttt{sos} and wait for
\texttt{>} and its bell. Enter one line:

% mips-batch: submit
\begin{commandexample}{Local SOS: submit the prepared jobfile}
batch test jobfile batchjob
\end{commandexample}

SOS asks \texttt{>password} and rings the bell. Only at that question enter:

% mips-batch: password
\begin{commandexample}{SOS password question: answer for the demo account}
test
\end{commandexample}

The dummy password can be visible; do not infer masked input or secure
authentication. The expected result, with a variable timestamp, is:

\begin{commandexample}{Expected local output: completed batch job}
08.00 sos: test enrolled
 batchprobe
08.00 sos: test removed after finis
\end{commandexample}

Enrollment alone is not completion. Wait for the actual message and removal
\emph{after finis}; \texttt{removed after break} is a failure, not the expected
result. The jobfile's \texttt{finis} terminates the batch FP automatically:
do not type FP commands at SOS while waiting. Select \texttt{s} and use
\texttt{list} if needed; the services remain present and \texttt{test} is gone.

After normal removal, the jobfile is unchanged and can be submitted again.
Request fresh attention and select SOS for the independently checked inline
alternative:

% mips-batch: inline
\begin{commandexample}{Local SOS: repeat with the inline dummy password}
batch test jobfile batchjob pass test
\end{commandexample}

This form does not ask a second password question. Expect the same single
enrollment, message and normal removal. Never resubmit merely because a reply
is delayed: first establish whether the previous job still exists.

\subsubsection{Check and remove the exercise input}\label{mips-ts-batch-cleanup}

Select \texttt{s} and create the same \texttt{bprep} FP again. Export the
existing guest file to a new host filename, then deliberately delete the
exercise entry:

\begin{operationnote}{Cleanup removes the retained jobfile}
Use these commands only for the \texttt{batchjob} created above and only after
both jobs have finished. They do not cancel a running job or remove host input
and output files. Preserve a backup instead if the guest file is wanted.
\end{operationnote}

% mips-batch: cleanup
\begin{commandexample}{Local FP bprep: export and remove the exercise jobfile}
batchcheck=tohost batchjob
clear user batchjob
lookup batchjob
finis
\end{commandexample}

Expect \texttt{***lookup batchjob unknown} and
\texttt{pause bprep finis fp}. The host export must have succeeded before
deletion; it refuses an existing \texttt{exchange/batchcheck}. Compare the
independent host copies with \texttt{cmp exchange/batchinput exchange/batchcheck}.
No output and status zero mean that execution preserved the jobfile bytes.

Historically, SOS schedules batch work in the background relative to
interactive \texttt{go}/\texttt{run} jobs. A jobfile supplies input while the
creating terminal receives output; this exercise submits from the local console.
It does not establish arbitrary output destinations or remote batch submission.
The plain \texttt{batch test} form does
not supply that jobfile and is not the checked recipe for unattended work.
Do not disconnect as a substitute for \texttt{finis}, or assume that this
single-account demonstration establishes concurrent-job operation.

\section{Files, persistence and shutdown}\label{mips-ts-persistence}
\index{MIPS/TS!persistence}

The demonstration account has a deliberately small disk claim: 30 temporary
segments/six entries at key 0, and retained-file claims of two entries/six
segments at each of keys 1 through 3. It is suitable for small FP experiments, not large compilation or
data workloads. Retained files still need \texttt{scope user}; temporary
files may be removed when the process finishes. See
\manualref{keep-guest-files-across-sessions}{Keep guest files across sessions}.
Do not assume an administrative FP's broader resource claims apply to this
account. The second disk primarily supplies SOS's permanent working space;
its existence does not automatically give the account that free space.

Unlike the default profile, this configuration retains writes automatically
in \texttt{UNIT0.mipsts.overlay} and \texttt{UNIT1.mipsts.overlay}. Both base
images remain read-only. Keep the two sidecars with their exact base images
and configuration. A complete restart with automatic subsystem startup and
exact readback of administrative and remote-created user-scope files has been checked; this
does not certify every application or file-recovery path.

\subsection{Keep a small file in the demo account}\label{mips-ts-retained-file}
\index{MIPS/TS!retained-file exercise}

This exercise uses the supplied unprotected profile, one remote \texttt{test}
session at a time and the existing small account quota. Work on private copies
with their own overlays, not historical or shipped base images in write-through
mode. The names \texttt{keepseed} and \texttt{keepsake} must be unused; do not
overwrite files from earlier work. Start the profile as above.

\subsubsection{Prepare the compact input}

First prepare a short host text. This subshell refuses to overwrite an existing
\texttt{exchange/keepinput}:

\begin{commandexample}[language=sh]{Host shell: create the exercise input}
(set -C; printf 'MIPS/TS file retained across sessions.\n' > exchange/keepinput)
\end{commandexample}

The file contains one line and a final LF:

% mips-files: input
\begin{commandexample}{Host file contents: exchange/keepinput}
MIPS/TS file retained across sessions.
\end{commandexample}

In the original terminal, press Escape and select \texttt{s}. Create a separate
administrative FP with catalog intervals matching the demo account:

% mips-files: administrator
\begin{commandexample}{Local Guest-s: create an administrative preparation FP}
all verify base 0 0 login 0 0 user 0 0 project 0 0 run
\end{commandexample}

Wait for \texttt{to verify} and the FP bell, then check both names:

% mips-files: absent
\begin{commandexample}{Local FP verify: check that the exercise names are unused}
lookup keepseed keepsake
\end{commandexample}

Expect \texttt{***lookup keepseed unknown} and
\texttt{***lookup keepsake unknown}. These diagnostics are intentional.
If either name is found, stop and preserve that file. Otherwise create a compact,
retained seed:

% mips-files: seed
\begin{commandexample}{Local FP verify: import and compact the seed}
keeptmp=fromhost keepinput
keepseed=copy keeptmp
scope user keepseed
finis
\end{commandexample}

Wait for \texttt{pause verify finis fp}. The large import belongs only to this
administrative process; \texttt{copy} makes the small seed fit in one segment.
The explicit 0/0 intervals make it visible to the demo job. This does not enlarge
the account's claims or create its own file on its behalf.

\subsubsection{Create the account's retained copy}

Connect by Telnet and log in with \texttt{run test pass test} as in
\manualref{mips-ts-login}{Connect and log in}. At the remote FP bell:

% mips-files: create
\begin{commandexample}{Remote FP test: create and retain your own copy}
keepsake=copy keepseed
scope user keepsake
lookup keepsake
c=copy keepsake
finis
\end{commandexample}

Expect a \texttt{keepsake =set 1 disc} descriptor ending in \texttt{; user},
the exact input line on the terminal and normal SOS removal after
\texttt{finis}. Plain \texttt{copy keepsake} reports copy statistics;
\texttt{c=copy keepsake} selects console output and displays the text.
The remote job creates and retains \texttt{keepsake}; the seed is only its input.

\subsubsection{Read it after login and restart}

Request attention, select SOS and log in again. Read the existing file without
recreating it:

% mips-files: read
\begin{commandexample}{Remote FP test: verify the retained file in a later session}
lookup keepsake
c=copy keepsake
finis
\end{commandexample}

The descriptor must still say \texttt{; user} and the line must be unchanged.

For byte-exact verification, select \texttt{s} locally and create the same
\texttt{verify} FP with 0/0 intervals. Before restarting use \texttt{keepcheck1};
after it use the new host export name \texttt{keepcheck2}:

% mips-files: export
\begin{commandexample}{Local FP verify: export the existing retained file}
lookup keepsake
keepcheck1=tohost keepsake
finis
\end{commandexample}

Both exports must be new host filenames; export refuses collisions. In a host
shell compare each with \texttt{exchange/keepinput}, for example
\texttt{cmp exchange/keepinput exchange/keepcheck1}. No output and exit status
zero mean equal bytes. A base-image listing does not include overlays and is
not a substitute for this live guest lookup/export.

After the first export, perform a complete normal stop and restart with the same
two overlays, following \manualref{mips-ts-normal-stop}{Stop, restart and back up}.
Log in again, repeat the read-only remote sequence, log out and export locally
as \texttt{keepcheck2}. Do not repeat seed preparation or
\texttt{keepsake=copy keepseed}: that would test recreation, not persistence.
Compare the second export with the host input before proceeding to cleanup.

\subsubsection{Remove only the exercise files}

\begin{operationnote}{Deliberate cleanup deletes the exercise files}
Only after checking the retained data, log in again and delete the remote
copy below. This removes that guest file; it is not a logout or overlay-reset
command. Keep a backup instead if its contents are wanted.
\end{operationnote}

% mips-files: cleanup
\begin{commandexample}{Remote FP test: remove only the exercise copy}
clear user keepsake
lookup keepsake
finis
\end{commandexample}

Expect \texttt{***lookup keepsake unknown}, followed by normal logout. Finally
create the local \texttt{verify} FP again, enter \texttt{clear user keepseed},
check that both names are unknown and finish with \texttt{finis}. Host input/
export files remain in \texttt{exchange/}; deleting a guest area does not delete
those independent host copies.

\subsection{Stop, restart and back up}\label{mips-ts-normal-stop}
\index{MIPS/TS!normal shutdown}\index{MIPS/TS!backup}

Before stopping, enter \texttt{finis} in the remote FP and wait for normal
SOS removal. Disconnect the Telnet client. In the original emulator terminal,
press Ctrl-E and enter \texttt{quit} at \texttt{sim4000>}. This exits the
emulator; it is not a guest system-wide shutdown command. Restart with the
same configuration and overlays to resume stored disk changes. The boot
date is set again on each start; it is not the current host date.

To return to the prepared demonstration, stop the emulator and remove both
MIPS/TS overlay files together. This loses all guest changes since the images
were prepared. Never remove an active overlay, reuse a different profile's
overlay, or run two emulators on the same pair. Do not run \texttt{kitlabel}
on the supplied second disk: it is already initialized.

For a backup, stop the emulator first and copy the matching
\texttt{UNIT0.mipsts}, \texttt{UNIT1.mipsts}, both overlay files,
\texttt{RC4000.mipsts.cfg} and \texttt{MONITOR} as one set. Preserve wanted
host files in \texttt{exchange/} separately. An offline catalog listing of
\texttt{UNIT0.mipsts} reads only that base image, not a running session's
sidecars; use guest \texttt{lookup} when checking current files.

\subsection{Reconnect after losing an idle FP connection}\label{mips-ts-idle-reconnect}
\index{MIPS/TS!Telnet reconnection}

This procedure is checked only for the supplied unprotected profile when FP
had finished its previous command and was waiting for input. Closing Telnet
without \texttt{finis} leaves that guest job running; it is not a logout.
Two successive disconnect/reconnect cycles preserve both a temporary file
and a retained user file. A TCP close is not a guest terminal hard error.

Keep the same emulator instance running. In its local terminal you can
inspect \texttt{test} with \texttt{list} in guest \texttt{s}; see
\manualref{mips-ts-status}{status checks}. At the simulator prompt,
\texttt{show telnet} should report device 4 as \texttt{listening} once the
old connection is closed. If you suspended the CPU with Ctrl-E, enter
\texttt{continue} before expecting remote guest replies.

Reconnect from a host terminal:

\begin{commandexample}[language=sh]{Host shell: reconnect to the same running instance}
telnet 127.0.0.1 4004
\end{commandexample}

Do not request attention or start another SOS job first. At the reconnected
terminal, send this directly to the already waiting FP:

\begin{commandexample}{Remote FP: confirm the existing session is responsive}
message resumed
\end{commandexample}

Expect the message \texttt{resumed} and the FP bell. No new enrollment or
password dialogue is expected. Use \texttt{lookup} on your known guest files
before continuing work; temporary files still belong to this same job.
Finish later with \texttt{finis} in the usual way. If the old job has already
ended or no FP reply arrives, stop and inspect its status rather than
replaying commands or removing the process blindly.

\begin{operationnote}{Reconnection is not a new login or a saved-session restore}
Anyone who can reach this local endpoint may resume the existing session
without its password. Keep the listener on loopback and restrict access to
trusted local users. Do not expose it as an authenticated remote service.

Lost output is not replayed. This clean-prompt recipe does not cover unknown
pending input, active transfers, programs still running, terminal hard errors or emulator
restart/crash recovery. Preserve wanted files and complete commands normally
before a planned disconnect; prefer \texttt{finis} when you intend to log out.
\end{operationnote}

\subsubsection{A known partial input line}\label{mips-ts-partial-input}
\index{MIPS/TS!partial input after reconnection}

Reconnecting does not clear characters already accepted by FP. The following
two harmless cases have been checked on the supplied profile. They require
that you know exactly what was typed and saw it echoed before the connection
was lost. They are not a method for guessing the contents of an interrupted
input buffer.

At a waiting \emph{remote FP}, type the following prefix \emph{without
pressing Enter}, then disconnect Telnet without \texttt{finis}:

% mips-reconnect: prefix
\begin{commandexample}{Remote FP: partial line, no Enter}
message par
\end{commandexample}

Reconnect to the same running emulator as above. Do not request attention or
log in again. Enter only the missing suffix and submit it:

% mips-reconnect: suffix
\begin{commandexample}{Reconnected remote FP: complete the existing line}
tial
\end{commandexample}

Expect \texttt{partial} followed by the FP bell: FP has executed
\texttt{message partial}, not a new command named \texttt{tial}.

For a correction example, start another line, again without Enter:

% mips-reconnect: correction-prefix
\begin{commandexample}{Remote FP: known partial line to correct, no Enter}
message wrong
\end{commandexample}

After disconnecting and reconnecting, press Ctrl-H exactly five times to
erase the known suffix \texttt{wrong}, then type \texttt{right} and Enter.
Ctrl-H is ASCII Backspace; the physical Backspace key is equivalent only
when its DEL output is mapped to Ctrl-H by the active configuration.
Expect \texttt{right} and the FP bell. This does not require removing the
job, discarding files or restarting the emulator.

Do not press Enter on an unknown partial line, and do not send a fresh whole
command on top of it. Bytes typed just before a network failure may not have
reached the guest; the examples only establish recovery of a fully echoed
known prefix. For anything uncertain, use
\manualref{mips-ts-unexpected-failure}{the non-destructive diagnostic steps}.

\section{Check status without disturbing the job}\label{mips-ts-status}

Keep the Telnet session for the user job. To inspect the guest from its
original local terminal, press Escape, select \texttt{s}, then enter:

\begin{commandexample}{Local guest system process s: inspect processes}
list
\end{commandexample}

The list should include \texttt{tem}, \texttt{primo} and \texttt{sos}.
During a login it also includes \texttt{test} under parent \texttt{sos};
after normal \texttt{finis}, that child disappears. A waiting interactive
child can be displayed as \texttt{stop} with stop count zero. That single
state word is not proof of a crash: successful FP input/output and normal
logout are the relevant checks. The resident services should remain active
and have zero stop counts. Process addresses and exact scheduling states
can change; do not copy addresses from an example into commands.

\texttt{list} here is an \texttt{s} command. Do not replace it with SOS
\texttt{display}: that input has returned \texttt{command unknown} in the
checked setup. Nor does a SOS \texttt{ready} response to a diagnostic
\texttt{call} prove login, password validation or a terminal connection.
Use the actual enrollment and FP exercise instead. Do not use \texttt{proc
test remove}, SOS \texttt{kill}, or a reset as routine logout commands;
they can discard the job's temporary work.

\section{Print through PRIMO}\label{mips-ts-primo-printing}
\index{PRIMO!printing}

The supplied configuration opens \texttt{printer.mipsts.txt} automatically.
No change to the disk images, \texttt{startprimo} or PRIMO generation is
required. This recipe uses the original local terminal for administration,
not the resource-limited \texttt{test} login. Start the MIPS/TS profile and
wait for \texttt{sos started} before proceeding.

In a second host terminal, prepare the two input files. Do not overwrite
existing exchange files without checking their contents first:

\begin{commandexample}{Host shell: prepare a program and printable text}
cp examples/primo_print.alg exchange/prbsource
printf 'Hello from PRIMO.\n' > exchange/prbinput
\end{commandexample}

In the original terminal, press Escape, enter \texttt{primo}, and wait for
its \texttt{=} prompt and bell. Sign up the real printer (kind 14):

\begin{commandexample}{Guest PRIMO operator: assign the physical printer}
signup printer 14
\end{commandexample}

Wait for the timestamped \texttt{primo:ready} reply. To inspect the
assignment, request attention again with Escape, select \texttt{primo},
then enter \texttt{display}. Expect a \texttt{printer} row ending in
\texttt{,idle}. Each operator command requires a fresh attention/selection;
the ready reply is not an invitation to type at a persistent prompt.

Press Escape and select \texttt{s}, then allocate an administrative FP:

\begin{commandexample}{Guest s: allocate a temporary printing session}
all prbprep size 150000 temp disc 500 20 run
\end{commandexample}

Wait for \texttt{to prbprep} and the FP bell. The following creates only
temporary guest entries in this process; do not substitute existing retained
files as output targets:

\begin{commandexample}{Guest FP: import, compile and submit one transport}
psrc=fromhost prbsource
ptmp=fromhost prbinput
pin=copy ptmp
plp=set lp printer
prbprog=algol psrc connect.no
prbprog
\end{commandexample}

Compilation must end with \texttt{algol end}, without errors. The program
calls the installed \texttt{transfer} procedure to define a transport, then
waits for it. Expect the following, with a positive, variable transport number:

\begin{commandexample}{Expected FP result: accepted and completed transport}
primo define 0 reply 0 transport 82
primo wait 0 reply 0 state 5 cause -1
end         26
\end{commandexample}

The numeric \texttt{end} value is also variable. State 5 with cause -1
denotes normal completion in this checked recipe. At that point the host
file already contains \texttt{Hello from PRIMO.} followed by LF; no emulator
shutdown is needed to flush it. The payload must not merely appear on the
console. \texttt{plp=set lp console1} would select a console substitute and
is deliberately not used here.

To print another text, replace the host \texttt{exchange/prbinput} with the
new input, then repeat \texttt{ptmp=fromhost prbinput}, \texttt{pin=copy ptmp}
and \texttt{prbprog} in the same FP. Do not change or remove \texttt{pin}
before the previous transport has completed. Keep lines within 132 columns;
the printer chapter describes character/control mapping and buffering limits.
The current PRIMO trimming emits the input without job banners or extra page
separators; do not assume that other PRIMO installations behave identically.

After completion, enter \texttt{finis} in FP and wait for
\texttt{pause prbprep finis fp}. Select PRIMO again and enter
\texttt{signoff printer}; wait for ready. A fresh \texttt{display} must no
longer show the assignment. Signoff removes the operator assignment; it does
not detach the emulator output file. Keep completed output when restarting:
the next run appends to it. The checked workflow includes repeated short and
160-line jobs, signoff, two boots with persistent disk overlays, and successful
TEM/SOS login/logout after printing on each boot.

This is a bounded administrative printing example, not a general SOS user
print command, concurrent queue test or protected
MIPS/TS acceptance. If an operation reports an error or waits for an operator,
retain its diagnostics and do not report the transport as completed.

\subsection{Recover a job submitted with output detached}
\label{mips-ts-printer-recovery}

When a transport is submitted while the emulator output is detached, the
checked PRIMO version reports \texttt{primo:***printer end document}.
This is an operator request, not successful completion. Select PRIMO and
enter \texttt{display}: the \texttt{printer} row identifies the input area
\texttt{pin} and ends in \texttt{,waiting}. Leave the submitting FP and its
input intact. Do not run \texttt{prbprog} again to retry the same transport.

Suspend with Ctrl-E. Reconnect the existing, healthy output file in append
mode, then resume CPU execution:

\begin{commandexample}{Emulator CLI: reconnect detached output}
attach 7 printer.mipsts.txt append
continue
\end{commandexample}

These paths are relative to the emulator's working directory. Select PRIMO
with Escape, then \texttt{primo}, and issue:

\begin{commandexample}{Guest PRIMO operator: resume the waiting transport}
start printer
\end{commandexample}

Wait for both \texttt{primo:ready} and the submitting program's completion:
\texttt{primo wait 0 reply 0 state 5 cause -1}, a normal FP \texttt{end}
and its input bell. The text is then present exactly once in the host file;
earlier output is preserved. Use \texttt{start}, not \texttt{restart}:
the latter deliberately starts the transport again from the beginning and
can repeat printed text. Do not use an exclusive attach to an existing file
or delete the output to make attachment succeed.

The acceptance test deliberately detaches output between completed jobs,
submits one new job, confirms the waiting state and unchanged file, then
reattaches and resumes. It also checks healthy TEM/PRIMO/SOS while waiting,
normal login/logout afterward and preservation across emulator restart.
This does not certify recovery from partial host writes, disk-full errors,
mid-job paper loss, unrestricted concurrent jobs or an emulator crash. Preserve diagnostics
and output rather than applying this recipe blindly to another failure.

\subsection{Multiple outstanding jobs: a bounded limit}

Two transports from separate administrative processes have been checked:
both are accepted while the printer is unavailable, then print exactly once
in submission order after reattachment and \texttt{start printer}. A further
transport completes on the same PRIMO instance after the second client is
removed. This does not establish arbitrary queue capacity or every scheduling
policy.

There is an important restriction: the second simultaneous blocking wait
returns reply 3 (\emph{missing resources}) in this session. The developer
fixture handles that reply explicitly by querying transport state with action
4 until completion, within finite test limits. It does not resubmit the data
or treat the refused wait as a successful print. The installed
\texttt{primotrim} generation text sets \texttt{waitops := 1}, unlike
\texttt{waittrans := 300}: queued transports and blocking waits have different
limits. A private test with eleven rather than ten runtime message buffers
still refuses the second wait. This strongly points to the generated PRIMO
waiter limit, not an emulator printer setting. Increasing startup buffers or
editing the generation text alone does not regenerate the installed program.
The PRIMO binary and shipped startup remain unchanged.

The supplied \texttt{primo\_print.alg} example is intentionally synchronous
and has no polling fallback. Use the sequential recipe above for ordinary
printing. If its wait reports a nonzero reply, preserve the transport number
and diagnostics: the original job may still be queued. Submitting it again
can create duplicate output. Limited test coverage is not a reason to enlarge
claims blindly or run \texttt{primotrim} on the installed system.

\section{Troubleshooting and current limits}\label{mips-ts-troubleshooting}
\index{MIPS/TS!troubleshooting}

First identify the context in \manualref{mips-ts-command-contexts}{the command
context guide}. Preserve the complete diagnostic and the selected config before
changing anything. The checks below do not authorize deleting wanted data or
resetting a working profile.

\begin{description}
\item[Connection refused.] Wait for startup, check that the emulator is still
running and that port 4004 is free. Check the selected configuration, not
the default \texttt{RC4000.cfg}, which has no active listener.
\item[No SOS prompt.] Press Ctrl-A in Telnet, then select \texttt{sos}.
If the local terminal is at \texttt{sim4000>}, enter \texttt{continue} there:
Ctrl-E suspends the whole machine, including remote work. A running FP normally
signals readiness with a bell, not a printed SOS or shell prompt.
\item[Bad password.] Use \texttt{run test pass test} at SOS, not at FP or the
emulator prompt. For the batch password question, answer \texttt{test} only
after \texttt{>password}; the inline batch form does not ask that question.
See \manualref{mips-ts-batch-submit}{batch submission}. Do not submit duplicate
jobs while diagnosing a delayed reply.
\item[Insufficient resources or an existing session.] Use one test job and
finish it normally. A tiny host import can still exceed the 30-segment temporary
claim; use the \manualref{mips-ts-retained-file}{compact-input procedure}, not
\texttt{kitlabel} or a larger disk. For a lost connection to an idle FP, use
\manualref{mips-ts-idle-reconnect}{the bounded reconnect procedure}; it does
not cover interrupted commands or a failed guest job. Extra accounts and
simultaneous terminals require deliberate guest provisioning, not just another
listener in the configuration.
\item[File unknown after login or restart.] Check whether it was retained with
\texttt{scope user}, whether the administrative FP used the specified 0/0
intervals, and whether both matching overlays were kept. A fresh overlay starts
from the supplied image, not the previous session's retained files. Inspect
with guest \texttt{lookup} before recreating or deleting anything.
\item[Only copy statistics appear.] Use \texttt{c=copy keepsake} to display the
exercise text. Plain \texttt{copy keepsake} is not the display command used in
the retained-file recipe.
\item[Host export refuses a file.] An existing export name must not be
overwritten. Preserve that host file and select a new name; for other export
failures, retain the diagnostic rather than assuming a name collision. Do not
run guest cleanup until the needed export succeeds.
\item[Batch enrolled but not finished.] Enrollment is only acceptance. The
checked jobfile must contain both \texttt{message batchprobe} and \texttt{finis}.
Wait for the message and \texttt{test removed after finis}; inspect \texttt{list}
from local \texttt{s} before deciding whether another submission is safe.
\item[Guest break or abnormal removal.] Record the originating process and
numbers. \texttt{removed after break} is not normal completion. Check memory/
storage claims before larger commands. Do not delete overlays containing wanted
work, force-remove processes or reset as an automatic recovery measure.
\item[Printer waiting or nonzero wait reply.] Preserve the transport number and
input area. Do not resubmit: the original job may still be queued. Follow
\manualref{mips-ts-printer-recovery}{detached-output recovery} only for that
specific checked failure, not disk-full or partial-write errors.
\item[Date prompt or missing disc1.] Select the matching \texttt{UNIT0.mipsts}
and \texttt{UNIT1.mipsts} pair. Older overlays can preserve previous startup
edits. Back up the stopped profile before resetting it; never initialize the
supplied second disk again.
\end{description}

\subsection{Preserve work after an unexpected failure}\label{mips-ts-unexpected-failure}
\index{MIPS/TS!failure diagnosis}

When none of the specific checked recipes applies, do not try a sequence of
SOS \texttt{stop}, \texttt{start}, \texttt{break} or \texttt{kill} commands.
Use the following order instead:

\begin{enumerate}
\item Stop sending remote input. Keep the complete terminal diagnostic and
the last command, noting whether Enter was sent, whether its output finished,
and which interpreter was receiving input. Do not replay an import, batch job
or print request merely because its acknowledgement was lost.
\item In the original terminal, suspend execution with Ctrl-E and inspect
the existing instance. These \emph{emulator CLI} commands do not remove a job:
\begin{commandexample}{Local emulator CLI: inspect before changing state}
show cpu
show telnet
show media
\end{commandexample}
Record the selected configuration and matching base/overlay names as well.
While suspended, the guest cannot answer FP or SOS commands. Enter
\texttt{continue} before a deliberate guest status check or export.
\item If the guest remains responsive, inspect processes through local
\texttt{s} as in \manualref{mips-ts-status}{the status procedure}. Inspect
wanted files with \texttt{lookup}; export them under new host names only
through a healthy FP with the appropriate intervals and resources. Successful
export, not a prompt alone, is the evidence that a wanted file was saved.
\item Prefer normal \texttt{finis} and shutdown only when the relevant jobs
are usable. If this is not possible, preserve the running instance and its
diagnostics before deciding to abandon it. Exiting loses core, pending input
and uncompleted work; an overlay is not a resumable machine snapshot.
\item After the emulator has stopped, preserve both overlay files with
their exact base images and configuration before further experiments.
Never copy active disks/sidecars as a consistent backup, reset by relabelling
\texttt{disc1}, or discard the only copy of wanted guest data.
\end{enumerate}

This is a preservation procedure, not a promise to recover every failed job.
Interrupted transfers/output, actual terminal hard errors, disk-full or
partial host writes, power loss and emulator crashes have no general checked
recovery sequence here. A future investigation must not be confused with a
supported instruction to mutate the current guest.

Punch/reader workflows and general printer error recovery are not certified.
The subsystem binaries are historical programs; their generation sources have
not been established as available on the image. Do not run \texttt{sostrim},
\texttt{temtrim} or \texttt{primotrim} as setup commands.

The default \texttt{UNIT0} does not start MIPS/TS. In \texttt{UNIT0.mipsts},
\texttt{startarea} reads the three component scripts directly; an extra
\texttt{read startmips} would exceed the monitor's two-level input stack and
report \texttt{not allowed}. \texttt{startsos} places its normal
2500-segment/200-entry claim on \texttt{disc1}, and \texttt{soscat} contains
the demo account. These differences are confined to the MIPS/TS images.

\section{Historical reference and coverage}\label{mips-ts-reference}

The historical \emph{Swopping Online System (SOS)}, RCSL 31-D662
(February 1982), explains job creation in section 2.2.1, printed pages
14--15; terminal access in section 2.3, pages 20--21; and user-catalog
resources in section 5.3.1. The \emph{MIPS/TS Operating Guide}, RCSL
31-D585, second edition (October 1979), supplies operator and subsystem
background. These manuals describe historical configurations, not a promise
that all their devices, accounts or commands are configured here. Scanned
OCR text can misread command syntax; prefer this chapter's checked recipes.

The ready-to-run coverage is one compatibility-monitor terminal session,
not multiterminal application attachment, invisible-password input,
arbitrary guest generation, general peripheral recovery, or protected
MIPS/TS. The historical SOS intervention commands \texttt{stop},
\texttt{start}, \texttt{break} and \texttt{kill} are not a certified recovery
sequence for this demonstration. In particular, historical \texttt{break}
behavior depends on system trimming and need not preserve a job's files.
